\section{Results}

\textit{This document is the \textbf{source of truth for all reported numbers.} Every value here traces to \texttt{RESEARCH\_LOG.md} \S{}B and the run directories under \texttt{trained\_models/}. Six results sections map to contributions C1--C6 plus external validity.}

\subsection{R1 --- Objective utility is conditional on the gap (C1)}

\textbf{Within domain}, objective choice is decisive. A linear probe on NTU features (30 epochs) spans \textasciitilde{}45 pp:

\begin{center}
\small
\begin{tabular}{ll}
\toprule
init & NTU linear-probe acc \\
\midrule
supervised (supLP120) & 59.6\% \\
hybrid (supMAE) & 59.2\% \\
contrastive (supcon) & 55.7\% \\
reconstruction (mae) & 22.6\% \\
scratch & 14.8\% \\
\bottomrule
\end{tabular}
\end{center}

Reconstruction alone (mae) barely clears random init within domain --- it does not build discriminative features. The two label-supervised objectives (supLP120, supMAE) and the label-supervised contrastive objective (supcon, Khosla 2020 SupCon --- in-batch same-label positives) all cluster in the mid-to-high 50s; pure reconstruction is far behind.

\textbf{Across the skeleton$\rightarrow$IMU gap}, the same choice compresses to a few points. LOSO-v2 final accuracy, per-(method,k) mean $\pm$ sd over 3 seeds (n=15):

\begin{center}
\small
\begin{tabular}{llll}
\toprule
method & k=0 & k=1 & k=3 \\
\midrule
supcon & 59.04 $\pm$ 6.83 & 82.62 $\pm$ 6.85 & 91.01 $\pm$ 3.96 \\
supMAE & 56.80 $\pm$ 6.70 & 82.57 $\pm$ 6.61 & 91.50 $\pm$ 4.24 \\
supLP120 & 57.67 $\pm$ 4.02 & 81.50 $\pm$ 6.34 & 91.25 $\pm$ 4.00 \\
scratch & 57.15 $\pm$ 6.84 & 82.31 $\pm$ 6.55 & 89.39 $\pm$ 4.90 \\
mae & 53.75 $\pm$ 5.07 & 79.16 $\pm$ 6.05 & 87.94 $\pm$ 4.37 \\
\bottomrule
\end{tabular}
\end{center}

The full method spread at k=1 is \textasciitilde{}3.5 pp (82.6 $-$ 79.2), vs the \textasciitilde{}45 pp within-domain spread. The gap, not the objective, sets how much objective choice matters. supcon leads narrowly at k=0/k=1 but the margin over scratch is not significant at any k (paired $\Delta$: +1.9/+0.3/+1.6 pp at k=0/1/3, all n.s., p$\ge$.12 --- see R2) --- a fifth objective, a fifth accuracy-is-a-wash story. All methods reach 94--98\% by k$\ge$5 (ceiling); k$\in$\{0,1,3\} is the discriminating regime.

\subsection{R2 --- A seed-stable negative-transfer case that inverts the field (C2)}

Pooled across 3 seeds (n=15 per method$\times$k cell), only one objective shows a significant, reproducible negative-transfer signature against scratch: pure reconstruction. What survives is not a supervised/contrastive collapse but a \textbf{reconstruction collapse} --- \texttt{mae} is the one objective consistently worse than scratch; everywhere else is a wash or a weak positive.

\textbf{Paired $\Delta$ vs scratch} (mean pp, [95\% CI], paired-t p), n=15:

\begin{center}
\small
\begin{tabular}{llll}
\toprule
contrast & k=0 & k=1 & k=3 \\
\midrule
supMAE $-$ scratch & $-$0.35 [$-$3.08, +2.39] p=.790 & +0.26 [$-$1.47, +2.00] p=.749 & \textbf{+2.11 [+0.88, +3.35] p=.003} \\
supLP120 $-$ scratch & +0.52 [$-$2.88, +3.93] p=.747 & $-$0.81 [$-$3.10, +1.49] p=.463 & +1.86 [$-$0.28, +4.00] p=.083 \\
supcon $-$ scratch & +1.90 [$-$2.56, +6.36] p=.377 & +0.31 [$-$3.21, +3.83] p=.853 & +1.62 [$-$0.49, +3.72] p=.122 \\
\textbf{mae $-$ scratch} & $-$3.39 [$-$6.91, +0.13] p=.058 & \textbf{$-$3.15 [$-$5.07, $-$1.23] p=.003} & $-$1.45 [$-$3.25, +0.35] p=.106 \\
\bottomrule
\end{tabular}
\end{center}

\textbf{Clip-level McNemar} (pooled 3 seeds; net = prior-only-correct $-$ scratch-only-correct):

\begin{table*}[!t]
\centering
\caption{R2 --- A seed-stable negative-transfer case that inverts the field (C2)}
\small
\begin{tabular}{lllllll}
\toprule
prior & k & n\_pairs & b (scratch-only) & c (prior-only) & net & p \\
\midrule
supMAE & 0 & 7974 & 536 & 514 & $-$22 & .517 (n.s.) \\
supMAE & 1 & 7644 & 327 & 349 & +22 & .419 (n.s.) \\
supMAE & 3 & 6984 & 162 & 312 & \textbf{+150} & 0.0 \\
supLP120 & 0 & 7974 & 731 & 778 & +47 & .236 (n.s.) \\
supLP120 & 1 & 7644 & 570 & 511 & $-$59 & .078 (n.s.) \\
supLP120 & 3 & 6984 & 250 & 385 & \textbf{+135} & 0.0 \\
supcon & 0 & 7974 & 715 & 877 & \textbf{+162} & 5e-05 \\
supcon & 1 & 7644 & 476 & 502 & +26 & .424 (n.s.) \\
supcon & 3 & 6984 & 272 & 390 & \textbf{+118} & 1e-05 \\
\textbf{mae} & 0 & 7974 & 737 & 471 & \textbf{$-$266} & 0.0 \\
\textbf{mae} & 1 & 7644 & 567 & 325 & \textbf{$-$242} & 0.0 \\
\textbf{mae} & 3 & 6984 & 379 & 280 & \textbf{$-$99} & .00013 \\
\bottomrule
\end{tabular}
\end{table*}

\textbf{mae is significant negative transfer at every k} (all p<.001) in both the pooled-accuracy and the clip-level view. It flips more clips wrong than right relative to random init, because it never encodes discriminative structure within-domain (R1) and the fine-tune from that init lands in a worse basin than from random. This inverts the field's working assumption that reconstruction/self-supervision is the \textit{safe} prior under large gaps.

The other three objectives are significant-positive \textbf{at k=3 only} (supMAE p=.003/McNemar p=0; supLP120 McNemar p=0; supcon both p<.001) and a wash at k=0/k=1 --- n.s. on both the paired-$\Delta$ and McNemar tests. The k=3 positive result for all three label-aware-or-hybrid objectives replicates cleanly across both tests; the k=0/k=1 "prior helps" claims for these three are smaller, less stable effects that should not be leaned on as heavily as the mae-negative-transfer claim.

\subsection{R3 --- Mechanism: gap is necessary-not-sufficient; CKA orders objectives (C3)}

\textbf{MMD does not track transfer.} Squared-MMD between NTU and Xsens-v2 features (supMAE encoder) is 0.0092 under v2 (below the local baseline 0.0109 and far below swing 0.0322), confirming v2 closes the gap --- but MMD ordering across encoders does not match transfer ordering (contrastive had near-lowest MMD yet worst transfer under the earlier preprocessing; mae has higher MMD than supervised yet similar low-k accuracy).

\textbf{CKA orders the objectives cleanly, but is a secondary check here --- the raw MMD\textsuperscript{2}/Frechet measurement below is the paper's primary domain-distance evidence, since it is the one that correctly characterizes objective performance against gap width.} Linear CKA between NTU and Xsens-v2 activations, per encoder $\times$ layer:

\begin{table*}[!t]
\centering
\caption{R3 --- Mechanism: gap is necessary-not-sufficient; CKA orders objectives (C3)}
\small
\begin{tabular}{lllll}
\toprule
encoder & proj & L0 & L1 & L2 \\
\midrule
supLP120 & 0.0038 & 0.0129 & 0.0134 & \textbf{0.0146} \\
supcon & 0.0038 & 0.0121 & 0.0114 & 0.0115 \\
supMAE & 0.0030 & 0.0044 & 0.0046 & 0.0044 \\
mae & 0.0026 & 0.0038 & 0.0034 & 0.0036 \\
scratch & 0.0026 & 0.0024 & 0.0024 & 0.0024 \\
\bottomrule
\end{tabular}
\end{table*}

Two facts, both important: (i) \textbf{absolute CKA is small for all} (<0.02) --- the cross-modal gap is large even after v2, so we do not over-claim on magnitude; (ii) the \textbf{ordering is clean and monotone with depth} for the two label-aware objectives --- supLP120 and supcon both align well above supMAE/mae/scratch, while scratch is flat. The best-aligned encoder is not the best-transferring on accuracy (supMAE is, R1), and the reconstruction encoder has higher CKA than scratch but negative transfer. \textbf{Conclusion:} representation alignment is a necessary context but not a sufficient predictor; the objective's discriminative inductive bias governs transfer. This is the mechanism the MMD-only view could not supply --- but the ordering-across-datasets claim rests on the raw, encoder-free measurement below, not on CKA.

\textbf{Extending this to the gap axis itself:} we also measured layer-wise CKA between NTU and each of the four targets used elsewhere in the paper (Xsens-v2, CZU skeleton, CZU IMU orientation quats, UTD skeleton) to test whether CKA orders them by our claimed small/middle/large gap (see R6e). It does not: on the supLP120 encoder, mean L0--L2 linear CKA is utd\_skeleton 0.0343 > czu\_skeleton 0.0253 > czu\_imu\_quat 0.0183 > xsens\_v2 0.0137 --- the claimed-middle target (xsens\_v2) ranks \textit{lowest}, not in the middle, and czu\_imu\_quat (claimed large gap) ranks above it. We report this as a further instance of necessary-not-sufficient rather than force a fit: CKA does not reliably track gap \textit{width} across independently collected datasets. This is exactly why the raw, encoder-free measurement below --- not CKA --- is the paper's evidence for the gap ordering. Full table in \texttt{wiki/concepts/domain-gap-metrics.md}.

\textbf{A raw, encoder-free measurement recovers the ordering CKA missed.} Both CKA and the encoder-space MMD above are computed \textit{inside a trained encoder's feature space} --- confounded by the pretraining objective, so they measure how different two datasets look \textit{to that encoder}, not how different the datasets are themselves. We therefore also measured domain gap directly on a hand-crafted, model-free feature (mean/std/var/skew/kurtosis over the 68 flattened LRQ channels --- the same feature used for the CRC published-baseline reproductions in R6/R6d, zero pretraining involved), computing squared-MMD (RBF kernel, median-heuristic bandwidth --- this fixes the fixed-sigma scale confound noted above) and Frechet distance (Gaussian closed form on a 50-component PCA) between NTU and each target, n=1000/side (n=861 for UTD, its full pool), seed 0:

\begin{center}
\small
\begin{tabular}{llll}
\toprule
target & gap (claimed) & MMD\textsuperscript{2} & Frechet \\
\midrule
czu\_skeleton & small & \textbf{0.0560} & \textbf{176.8} \\
utd\_skeleton & small & \textbf{0.0960} & \textbf{202.1} \\
xsens\_v2 & middle & \textbf{0.1218} & \textbf{281.9} \\
czu\_imu\_quat & large & \textbf{0.2049} & \textbf{393.8} \\
\bottomrule
\end{tabular}
\end{center}

\textbf{Both metrics confirm the claimed small<small<middle<large gap ordering exactly and monotonically} --- the opposite outcome from the CKA-per-target result immediately above. Removing the pretrained encoder from the measurement, rather than adding one, is what recovers the ordering: CKA measures how well a \textit{specific trained encoder} aligns two domains (task- and objective-dependent, hence the necessary-not-sufficient finding above), while raw MMD/Frechet measure how different the \textit{datasets themselves} are before any model touches them. The two results are complementary, not contradictory: the objective's discriminative inductive bias governs whether a given alignment \textit{transfers} (CKA's lesson), while the raw distributional gap between datasets follows the intuitive device/modality ordering (this result) --- CKA's failure to track it is itself evidence that encoder alignment and raw distributional distance are different quantities. This gives the R6e gap ordering a measured, encoder-independent basis it previously lacked. Full table in \texttt{wiki/concepts/domain-gap-metrics.md}, script \texttt{scripts/main\_experiment/raw\_domain\_gap.py}.

\subsection{R4 --- The prior's real value: calibration and data-efficiency (C4$\rightarrow$A2, C5)}

The accuracy spread is compressed, but the prior still pays off in three deployment-relevant ways.

\textbf{(a) Calibration.} Temperature-scaled ECE, per method$\times$k:

\begin{center}
\small
\begin{tabular}{llll}
\toprule
method & k=0 & k=1 & k=3 \\
\midrule
supLP120 & \textbf{0.0295} & \textbf{0.0282} & \textbf{0.0291} \\
supcon & 0.0298 & 0.0415 & 0.0318 \\
scratch & 0.0676 & 0.0494 & 0.0396 \\
supMAE & 0.0602 & 0.0520 & 0.0498 \\
mae & 0.0728 & 0.0577 & 0.0560 \\
\bottomrule
\end{tabular}
\end{center}

supLP120 is best-calibrated at every k; supcon is a clear second, ahead of supMAE/scratch/mae at every k; mae is worst at every k. The two label-supervised objectives (supLP120, supcon) lead calibration; the two without direct softmax-style label supervision at the pooled embedding (supMAE, mae) trail --- calibration tracks label-awareness of the objective, not just accuracy. This ranking is the most stable result in the paper. (Raw ECE at k=0 is large for all, 0.26--0.33, T$\approx$2.7--3.4 --- everyone is overconfident before temperature scaling.)

\textbf{(b) Calibration efficiency (AUC-30 = mean eval-acc over the first 30 calibration epochs) and convergence speed.} Paired $\Delta$ vs scratch:

\begin{itemize}
  \item AUC-30: supMAE k3 \textbf{+2.04 [+0.26, +3.81] p=.028}; supLP120 k3 \textbf{+3.74 [+1.03, +6.46] p=.010}; supcon k3 \textbf{+4.18 [+1.44, +6.93] p=.006} (k1 +1.89, n.s.); mae k1 $-$4.37 p<.001, k3 $-$3.83 p<.001.
  \item Convergence (first epoch $\ge$90\% of final; lower=faster): supLP120 k1 \textbf{$-$4.00 [$-$6.97, $-$1.03] p=.012}, k3 \textbf{$-$3.40 [$-$5.07, $-$1.73] p=.001}; supcon k3 \textbf{$-$3.53 [$-$5.64, $-$1.43] p=.003} (k1 $-$2.47, n.s.); mae \textit{slower} (+2.93 p$\le$.034); supMAE $\approx$ scratch.
\end{itemize}

The prior converges faster and integrates the few calibration shots more efficiently even where the endpoint accuracy is a wash --- and this holds for \textbf{two} label-aware objectives (supLP120, supcon), not one, strengthening R4 from a single-objective observation to a pattern tied to label-awareness rather than one specific pretraining recipe.

\textbf{(c) Subject-count scaling (A2) --- how long the prior stays useful.} Prior benefit vs scratch (pp), by number of fine-tuning subjects N $\times$ k, \textbf{pooled over 3 seeds (42/43/44), n=15 per cell}:

\begin{table*}[!t]
\centering
\caption{R4 --- The prior's real value: calibration and data-efficiency (C4$\rightarrow$A2, C5)}
\small
\begin{tabular}{lllllll}
\toprule
prior & k & N=0 & N=1 & N=2 & N=3 & N=4 \\
\midrule
supLP120 & 1 & +7.35 & \textbf{+18.20} & +5.02 & +1.14 & $-$0.81 \\
supLP120 & 3 & +18.22 & \textbf{+28.25} & +7.57 & +2.66 & +1.86 \\
supMAE & 1 & +3.15 & +5.20 & +2.25 & $-$0.34 & +0.26 \\
supMAE & 3 & +8.58 & +8.92 & +1.95 & +0.84 & +2.11 \\
mae & 1 & +1.35 & $-$0.42 & $-$1.74 & $-$2.94 & $-$3.15 \\
mae & 3 & +4.39 & +3.52 & $-$1.84 & $-$0.89 & $-$1.45 \\
\bottomrule
\end{tabular}
\end{table*}

Paired-t (n=15): supLP120 $-$ scratch is significant at N=0/N=1 both k (p$\le$.008, most p<.0001); supMAE $-$ scratch likewise significant at N=0/N=1 both k (p$\le$.008); mae $-$ scratch turns significantly negative at N=2 k=0 (p=.007), N=3 k=0/k=1 (p$\le$.02), and N=4 k=1 (p=.003) --- n.s. at N=1.

The prior's benefit \textbf{peaks at N=0/1} (supLP120 \textbf{+28.3 pp @ k=3, N=1}) and \textbf{washes to $\approx$1--2 pp by N=4} --- a clean, statistically confirmed deployment lever: the pretrained prior is worth roughly a full retraining set when 0--1 users are enrolled and is nearly redundant once \textasciitilde{}4 are. mae is negative from N$\ge$2 onward --- reinforcing R2.

\textbf{Caveat established in R6c below: this cold-start lever is itself target-representation-contingent.} Repeating this exact sweep on CZU's strong (raw-signal) target finds no N at which either prior beats scratch --- both are significantly \textit{worse} at N=0, k=1. Read R4c and R6c's cold-start extension together, not R4c alone.

\subsection{R5 --- From recognition to control (C6)}

The abstract controller (12-step Sequential Control Task, held-out posteriors, confidence-reject safety layer, asymmetric cost) turns the compressed accuracy differences into task-level ones. The System Input assignment (two Safety-Critical States, five Routine States) is hypothetical, not a recorded gesture set --- real gestures are relabeled onto it, and deliberately not by any property that could be seen as tuned to a preferred outcome (Methods \S{}8). A controller has design knobs (System Input assignment, error-cost model, operating threshold); rather than freeze any of them at one defensible-looking value, we probe all three simultaneously with a single Monte Carlo protocol: \textbf{120 independently, uniformly-randomly drawn 7-gesture System Input assignments}, each evaluated with \textbf{1000 Monte Carlo task-execution trials per (assignment, method, k, $\tau$/C\_crit) cell}. No assignment is chosen by any property of the recognizer (e.g. per-gesture recall) --- every assignment is a fresh unweighted random draw of 7 of the 22 recorded gestures, so no lock's finding can be attributed to a favorably- or unfavorably-picked task design. Three locks apply this same randomized-assignment protocol to three different design knobs:

\textbf{Lock 1 --- randomized assignment, base outcome model.} Mean hard task-success across the 120 assignments (1000 trials/assignment/cell):

\begin{table*}[!t]
\centering
\caption{R5 --- From recognition to control (C6)}
\small
\begin{tabular}{llllll}
\toprule
condition & scratch & mae & supMAE & supLP120 & supcon \\
\midrule
k=1, $\tau$=0 (ungated) & 0.516 & \textbf{0.455} & 0.522 & 0.472 & 0.523 \\
k=1, $\tau$=0.9 (gated) & 0.714 & \textbf{0.650} & 0.717 & 0.703 & 0.740 \\
k=3, $\tau$=0 & 0.663 & \textbf{0.643} & 0.722 & 0.723 & 0.713 \\
k=3, $\tau$=0.9 & 0.830 & \textbf{0.792} & 0.861 & 0.883 & \textbf{0.895} \\
\bottomrule
\end{tabular}
\end{table*}

\textbf{mae is the worst-compounding init in all four conditions.} supMAE and supcon are both $\ge$ scratch in every condition; supcon is the single best method at k=3, $\tau$=0.9.

\textbf{Lock 2 --- critical-cost severity, same 120 randomized assignments.} Replacing binary task-failure with a recoverable critical penalty C\_crit and sweeping C\_crit $\in$ \{2, 5, 10, 20, 50, $\infty$\}, k=1, median mean task-cost across the 120 assignments (1000 trials/assignment/cell):

\begin{table*}[!t]
\centering
\caption{R5 --- From recognition to control (C6)}
\small
\begin{tabular}{llllll}
\toprule
C\_crit & mae & scratch & supLP120 & supMAE & supcon \\
\midrule
20 & \textbf{32.3} & 29.0 & 30.4 & 28.5 & 28.6 \\
50 & \textbf{55.5} & 48.3 & 51.2 & 47.1 & 47.8 \\
$\infty$ & \textbf{756,017} & 644,016 & 704,016 & 611,518 & 638,517 \\
\bottomrule
\end{tabular}
\end{table*}

\textbf{mae has the highest median mean-cost at every C\_crit swept, confirming Lock 1's ordering under a completely different (and much harsher) outcome model.} supLP120 trends second-worst as the penalty grows, but never overtakes mae; supMAE is consistently the cheapest method.

\textbf{Lock 3 --- iso-safety operating point, same 120 randomized assignments.} Rather than pick $\tau$, we fix a false-activation \textit{budget} per assignment and find the smallest $\tau$ meeting it per method, then report the distribution across the 120 assignments. At a \textbf{1\% budget}:

\begin{table*}[!t]
\centering
\caption{R5 --- From recognition to control (C6)}
\small
\begin{tabular}{lllll}
\toprule
k & method & $\tau$* (median) & mean task-success & mean cost \\
\midrule
1 & \textbf{mae} & 0.97 & \textbf{0.360} & 22.98 \\
1 & scratch & 0.97 & 0.514 & 23.14 \\
1 & supLP120 & 0.99 & 0.561 & 24.18 \\
1 & supMAE & 0.99 & 0.524 & 24.30 \\
1 & supcon & 0.97 & 0.646 & 23.22 \\
3 & \textbf{mae} & 0.90 & \textbf{0.741} & 21.77 \\
3 & scratch & 0.90 & 0.797 & 20.80 \\
3 & supLP120 & 0.85 & \textbf{0.862} & 18.41 \\
3 & supMAE & 0.90 & 0.812 & 19.47 \\
3 & supcon & 0.85 & 0.848 & 18.34 \\
\bottomrule
\end{tabular}
\end{table*}

\textbf{mae has the lowest mean task-success at the 1\% budget at both k=1 and k=3} --- the worst method under Lock 3 too, at both shot counts and both budgets tested (1\% and the stricter 0.5\%, not tabulated). supLP120 and supcon lead at k=3, consistent with them being the two best-calibrated objectives (R4a).

\textbf{Net: all three locks agree.} Across 120 randomized System Input assignments, both outcome models (hard task-success and soft cost), a 6-point critical-cost sweep, and a tuning-free iso-safety operating point --- \textbf{mae compounds worst, consistently, under every stress test.} One nuance persists across the randomization: supLP120 still dips just below scratch on the \textit{ungated} Lock 1 metric at k=1 (0.472 vs 0.516) and trends as the second-most-expensive method as Lock 2's penalty grows severe --- both consistent with a specific, previously-documented mechanism (supLP120 occasionally maps an unrelated gesture onto a highly-separable Safety-Critical-State anchor with high confidence). This does not contradict the calibration story (R4a: supLP120 remains the best-calibrated method on average, and leads at k=3 under iso-safety) --- it is a secondary, narrower effect, visible only at the ungated operating point or under severe penalties, that never rises to displace mae as the worst-compounding objective under any of the three locks.

\subsection{R6 --- External validity (CZU-MHAD skeleton$\rightarrow$skeleton)}

Same objectives and LOSO protocol on CZU-MHAD's skeleton modality (5 subjects, 22 actions). Final accuracy, mean over 5 subjects (single seed, to match the CRC same-splits head-to-head below; pooled 3-seed results follow):

\begin{table*}[!t]
\centering
\caption{R6 --- External validity (CZU-MHAD skeleton$\rightarrow$skeleton)}
\small
\begin{tabular}{lllll}
\toprule
k & scratch & mae & supMAE & supLP120 \\
\midrule
0 (zero-shot) & 84.9 & 85.0 & 89.0 & \textbf{89.9} \\
1 & 90.2 & 88.5 & 92.8 & \textbf{92.5} \\
3 & 91.9 & 91.5 & 94.9 & \textbf{95.0} \\
\bottomrule
\end{tabular}
\end{table*}

Paired $\Delta$ vs scratch (seed 42): supLP120 and supMAE both clearly positive at zero-shot; \textbf{mae is essentially flat here, not negative} (+0.1 @ k=0), unlike the significant negative transfer it shows on NTU$\rightarrow$Xsens IMU. The mae-hurts signature does \textbf{not} reproduce on a same-modality (skeleton$\rightarrow$skeleton, small-gap) target --- evidence that the negative transfer is specific to the \textit{cross-modal} gap, which is exactly the mechanism R2/R3 predict. This is skeleton$\rightarrow$skeleton, not the IMU gap; the true cross-modal replication on CZU's inertial modality is reported next (R6b).

\textbf{Multi-seed robustness (3 seeds 42/43/44, pooled n=15).} The zero-shot supLP120 win holds: \textbf{supLP120 $-$ scratch = +5.00 pp at k=0, significant (paired-t p=.0012)}. supMAE is significant-positive at k=0 (+3.15 pp, p=.010) and k=1 (+2.11 pp, p=.032); its k=3 margin is a trend (+1.85 pp, p=.086). mae is positive-but-n.s. at k=0, slightly negative-but-n.s. at k=1/k=3 --- the same-modality "reconstruction does not clearly hurt" signature holds directionally. The small-gap "supervised prior wins" claim is seed-stable, and is independently replicated on a second public dataset (UTD-MHAD, R6d).

\textbf{supcon (3 seeds, pooled n=15) wins even bigger.} supcon $-$ scratch = \textbf{+5.24 / +4.67 / +3.41 pp at k=0/1/3, all significant (p=.0015 / .0003 / .0049)} --- a larger and more consistent margin than supLP120's. supcon > supLP120 in 28/42 folds (sign p=.044, mean$\Delta$ +0.91 pp): on this dataset the \textit{contrastive} label-aware objective edges out the \textit{softmax} label-aware objective, but both dominate scratch by a wide margin. The small-gap story is therefore not "supervised classification specifically wins" but \textbf{"any objective with direct label supervision at the pooled-embedding level wins when the gap is small."}

\textbf{Published-baseline comparison.} The CZU-MHAD dataset paper (Chao et al., 2022) reports a Collaborative Representation Classifier (CRC) on statistical-moment features; its protocol-comparable \textit{cross-subject} skeleton accuracy is \textasciitilde{}75.5\% (its closed, subject-mixed test is not comparable to LOSO). We reproduce that baseline family --- mean/std/var/skew/kurtosis features + CRC ($\lambda$=1e-4) --- on the \textbf{byte-identical LOSO k-shot splits} used by our learned recognizer, giving a same-splits head-to-head (mean acc over 5 folds):

\begin{table*}[!t]
\centering
\caption{R6 --- External validity (CZU-MHAD skeleton$\rightarrow$skeleton)}
\small
\begin{tabular}{llllll}
\toprule
k & CRC baseline & scratch & mae & supMAE & supLP120 \\
\midrule
0 & 84.81 & 84.87 & 84.95 & 89.03 & \textbf{89.94} \\
1 & 89.31 & 90.15 & 88.48 & 92.76 & \textbf{92.54} \\
3 & 90.82 & 91.90 & 91.46 & 94.92 & \textbf{95.04} \\
\bottomrule
\end{tabular}
\end{table*}

Two things stand out. (i) \textbf{The from-scratch learned encoder is statistically level with the hand-crafted CRC baseline} (84.87 vs 84.81 at k=0) --- architecture alone is not the advantage. (ii) \textbf{Only the NTU-pretrained prior opens a clear, consistent margin}: supLP120 $-$ CRC = +5.1 / +3.2 / +4.2 pp at k=0/1/3. On an independent public dataset the value is carried by the \textit{transferred prior}, not the recognizer, corroborating R4. (Our CRC reproduction, 84.8\%, exceeds the paper's published cross-subject skeleton \textasciitilde{}75.5\% --- expected given our LRQ representation and 4-subject LOSO dictionary vs their raw-position T5--T7 splits; we report it as a same-splits reproduction, not a claim about their exact pipeline.)

\subsubsection{R6b --- Cross-modal external validity (CZU-MHAD skeleton$\rightarrow$inertial)}

The true cross-modal replication: NTU skeleton prior $\rightarrow$ CZU \textbf{wearable-IMU} target (10 body-worn 6-axis sensors, \textbf{no magnetometer}), an independent public instance of the skeleton$\rightarrow$inertial gap. To feed the skeleton-pretrained encoder, the raw inertial is reduced to 17-segment orientation quaternions via Madgwick AHRS (yaw drifts without magnetometer) --- an orientation-only representation shared by every method. Final accuracy, mean over LOSO folds (pooled 3 seeds 42/43/44 $\times$ 5 subjects = 45 folds; CRC column single-seed, raw accel+gyro):

\begin{table*}[!t]
\centering
\caption{R6b --- Cross-modal external validity (CZU-MHAD skeleton$\rightarrow$inertial)}
\small
\begin{tabular}{llllll}
\toprule
k & scratch & mae & supMAE & supLP120 & CRC (raw accel+gyro) \\
\midrule
0 & 55.7 & 56.2 & \textbf{56.4} & 54.9 & 86.8 \\
1 & 54.3 & 54.8 & 54.5 & \textbf{52.8} & 90.8 \\
3 & 58.3 & 58.1 & \textbf{59.6} & 57.6 & 94.4 \\
\bottomrule
\end{tabular}
\end{table*}

\textbf{Under an impoverished (orientation-only) target encoding, the supervised prior's advantage disappears and trends toward hurting.} On the same-modality skeleton target (R6) the pure-supervised prior supLP120 dominated (+5.0 pp zero-shot, p=.001); across the modality boundary \textbf{supLP120 trends worst --- below scratch at every k}, though not significantly so at this sample size (supLP120 $-$ scratch = $-$1.0 pp mean; paired-t k=0 $-$0.86 pp p=.468, k=1 $-$1.44 pp p=.412, k=3 $-$0.70 pp p=.592). The \textbf{prior-vs-prior contrast (supMAE > supLP120) points the same direction}: 29/43 folds (sign p=.032, mean$\Delta$ +1.74 pp). supMAE does not clearly beat scratch (sign test 28/43 folds, p=.066, mean$\Delta$ +0.74 pp; paired-t clears nothing at any k). The honest reading: \textbf{supMAE $\approx$ scratch $\approx$ supLP120 at this gap once significance is accounted for} --- the \textit{direction} every prior points (reconstruction-containing objectives hold up better than the pure-supervised one) is mechanistically consistent with R2/R3, but this section alone does not carry a clean statistically-significant inversion. R6c below (dual-branch, on a stronger target) is where the sharper, significant version of this story lives.

Two honest limits. (i) The deep encoders trail the raw-inertial CRC baseline in absolute terms (55--60\% vs 87--94\%): the orientation-only encoding needed for cross-modal transfer discards the accelerometer-magnitude signal CRC exploits and carries yaw drift, so these numbers index \textit{relative} prior transfer, not the inertial ceiling. (ii) Unlike NTU$\rightarrow$Xsens, mae here is not a negative-transfer culprit (it trends slightly \textit{positive} vs scratch, +0.74/+0.51/$-$0.26 pp at k=0/1/3, all n.s.) --- mae's negative transfer is cross-domain-specific, not cross-modal-specific. Treat the direction in this section as the load-bearing claim, and the specific p-values as limited by sample size (n=43--45).

\textbf{supcon (3 seeds, pooled n=15/k) tracks supLP120's negative-to-neutral pattern.} supcon $-$ scratch = \textbf{$-$0.44 / $-$1.00 / $-$1.04 pp @ k=0/1/3} (all n.s., p$\ge$.55); pooled sign test supcon > scratch is 24/44 (p=.65, mean$\Delta$ $-$0.83, a wash trending negative). supcon vs supLP120 is statistically even (21/43, p=1.0, mean$\Delta$ +0.17) --- the two label-aware objectives land in the same negative-to-neutral band, both numerically below supMAE. \textbf{Objectives lacking a reconstruction/regularization component do not clearly help, and trend toward hurting, across this modality gap}, regardless of whether label supervision is delivered via cross-entropy (supLP120) or contrastive alignment (supcon) --- a directional pattern corroborated by R6c's sharper, significant version of the same contrast.

\subsubsection{R6c --- The prior's value is contingent on target-representation poverty (CZU-MHAD dual-branch)}

R6b's orientation-only encoding capped deep accuracy at 56--61\% --- far below the raw-inertial CRC baseline. Two questions follow: was that the \textit{encoding} or the \textit{architecture}, and does the R6b reconstruction-prior benefit survive once the target model uses its \textbf{full raw signal}? We test both with a dual-branch recognizer (a from-scratch branch on the raw 10$\times$6 accel+gyro stream + the NTU-pretrained encoder on the R6b orientation quaternions, concatenated to a shared head) on the byte-identical CZU LOSO splits, pooled over 3 seeds $\times$ 5 subjects (45 folds):

\begin{table*}[!t]
\centering
\caption{R6c --- The prior's value is contingent on target-representation poverty (CZU-MHAD dual-branch)}
\small
\begin{tabular}{lllllll}
\toprule
mode & scratch & mae & supMAE & supLP120 & supcon & CRC \\
\midrule
raw-only (k=0/1/3) & 81.9 / 90.0 / \textbf{95.6} & --- & --- & --- & --- & 86.8 / 90.8 / 94.4 \\
dual (k=0) & \textbf{86.0} & 85.3 & 85.9 & 84.5 & 83.8 & --- \\
dual (k=1) & \textbf{88.8} & 88.0 & 88.5 & 87.2 & 86.9 & --- \\
dual (k=3) & \textbf{91.9} & 91.4 & 91.6 & 90.9 & 90.0 & --- \\
\bottomrule
\end{tabular}
\end{table*}

Two results. \textbf{(i) Representation, not architecture, was the R6b bottleneck.} The raw-signal branch alone reaches 81.9/90.0/95.6, matching/beating CRC and \textasciitilde{}30 pp above R6b's orientation-only encoders --- the collapse there was the lossy encoding, not the deep model. \textbf{(ii) On a strong target the NTU prior adds no value, and label-supervised-only priors actively hurt.} \texttt{dual/scratch} is the best performer at every k, beating every prior: supMAE ties it (pooled 18/37 folds, sign p=1.0 --- no detectable difference), mae trends below (15/42, p=.088), \textbf{supLP120 is significantly worse --- pooled 10/41 folds (sign p=.0015), paired-t significant at k=0} ($-$1.50 pp, p=.017), and \textbf{supcon is worse still --- the single worst pooled result in this dataset (3/42 folds, sign p<.0001)}, paired-t significant at every k ($-$2.17 / $-$1.90 / $-$1.89 pp, p=.0002 / .0002 / .0066). The reconstruction component, not the specific supervised recipe, is what keeps an objective from actively hurting at large gap; the two label-aware-only priors (supLP120, supcon) are the clearest losers in the whole paper on this target.

\textbf{Dose-response along target richness.} Interpolating a middle rung between R6b (orientation-only) and R6c (full raw) --- a 20-channel per-sensor accel-magnitude + gyro-magnitude target --- traces the prior benefit as a smooth function of target strength ($\Delta$ vs scratch, k=0/1/3, pooled seeds 42/43/44):

\begin{center}
\small
\begin{tabular}{llll}
\toprule
target representation & scratch acc & supMAE $\Delta$ & supLP120 $\Delta$ \\
\midrule
0-ch quat-only (R6b) & 56 / 54 / 58 & +0.7 / +0.2 / +1.3 (n.s.) & $-$0.9 / $-$1.4 / $-$0.7 (n.s.) \\
20-ch magnitudes (dial) & 84 / 88 / 90 & $-$1.3 / $-$1.2 / $-$0.0 & $-$3.8 / $-$5.7 / $-$3.2 \\
60-ch full raw (R6c) & 86 / 89 / 92 & $-$0.1 / $-$0.3 / $-$0.4 & $-$1.5 / $-$1.6 / $-$1.1 \\
\bottomrule
\end{tabular}
\end{center}

The reconstruction prior's edge, where it exists at all, is concentrated at the fully-crippled orientation-only target and fades as the target sees more real motion; supLP120 is negative at every rung and significantly so from the 20-ch dial onward.

\textbf{Does the cold-start advantage (R4c) survive on a strong target?} R4c's deployment claim --- the prior is worth most at 0--1 enrolled subjects --- was established only on Xsens, a weak/moderate target. We repeat the subject-count sweep (N=0..3, \texttt{--n-train-subjects} added to \texttt{scripts/external/czu/dualbranch.py}) on CZU-dual, single-seed (n=5 folds; a 3-seed extension is in progress):

\begin{table*}[!t]
\centering
\caption{R6c --- The prior's value is contingent on target-representation poverty (CZU-MHAD dual-branch)}
\small
\begin{tabular}{lllll}
\toprule
N & k & scratch & supLP120 $\Delta$ & supMAE $\Delta$ \\
\midrule
0 & 0 & 2.87 & +0.11 (n.s.) & +0.37 (n.s.) \\
0 & 1 & 38.94 & \textbf{$-$5.62 p=.045} & \textbf{$-$5.55 p=.005} \\
0 & 3 & 54.21 & +0.23 (n.s.) & \textbf{$-$2.87 p=.035} \\
1 & 1 & 71.01 & $-$3.11 p=.085 (n.s.) & $-$3.51 p=.052 (n.s.) \\
1 & 3 & 79.58 & \textbf{$-$2.54 p=.004} & $-$1.80 p=.076 (n.s.) \\
\bottomrule
\end{tabular}
\end{table*}

\textbf{Neither prior beats scratch at any N, and both are significantly worse at N=0, k=1} --- the exact cold-start cell where the Xsens A2 result showed the prior's \textit{largest} benefit (+18.4 pp, R4c). This scopes the deployment claim precisely: \textbf{the cold-start advantage is itself contingent on target representation poverty}, mirroring R6c's main finding at N=4. A strong target's from-scratch branch reaches a useful operating point with zero enrolled subjects and no prior --- there is no regime tested here where the NTU prior helps on this target.

Read across the three CZU settings, the prior's value \textit{and its ranking} degrade monotonically with the gap:

\begin{table*}[!t]
\centering
\caption{R6c --- The prior's value is contingent on target-representation poverty (CZU-MHAD dual-branch)}
\small
\begin{tabular}{lllll}
\toprule
Setting & Gap & Target repr. & Best prior & supLP120 (pure-supervised) \\
\midrule
R6 skeleton$\rightarrow$skeleton & small (same modality) & strong (native skeleton) & \textbf{supLP120 +5.0 pp @k0} (p=.001), beats CRC & best-in-class \\
R6b IMU, orientation-only & large (cross-modal) & weak (orientation-only) & none clearly (supMAE trends > supLP120, 29/43, p=.032) & trends worst, n.s. (p=.41--.59) \\
R6c IMU, raw/dual & large (cross-modal) & strong (raw $\approx$ CRC) & none (\texttt{dual/scratch} itself is best at every k) & worst, p=.0015 \\
\bottomrule
\end{tabular}
\end{table*}

The two axes are separable: holding the gap large, moving the target representation from weak (R6b) to strong (R6c) removes the reconstruction prior's advantage; holding the target strong, moving the gap from small (R6, where supervised wins) to large (R6c) turns the supervised prior from best to worst. A single public dataset yields a controlled contrast in which the \textit{same} supervised prior is best-in-class in one modality and worst-in-class in another, purely as a function of gap width and target capability --- the sharpest external corroboration of the negative-transfer thesis (R2/R3): a transferred prior earns its keep only in proportion to the target's representational poverty.

\subsubsection{R6d --- Second independent same-modality dataset (UTD-MHAD skeleton$\rightarrow$skeleton)}

To confirm the small-gap "supervised prior wins" prediction (R6) is not specific to CZU, we replicate it on a second, independently collected public dataset: UTD-MHAD (Kinect-v1, 20 joints remapped to the NTU-25 layout; 8 subjects, 27 actions, 861 clips). Same objectives and LOSO k-shot protocol; pooled over 3 seeds $\times$ 8 subjects (24 folds per cell):

\begin{table*}[!t]
\centering
\caption{R6d --- Second independent same-modality dataset (UTD-MHAD skeleton$\rightarrow$skeleton)}
\small
\begin{tabular}{lllll}
\toprule
k & scratch & mae & supMAE & supLP120 \\
\midrule
0 (zero-shot) & 77.4 & 79.6 & 80.0 & \textbf{82.4} \\
1 & 90.9 & 91.6 & 93.1 & \textbf{93.0} \\
3 & 94.0 & 94.5 & 95.1 & \textbf{96.0} \\
\bottomrule
\end{tabular}
\end{table*}

\textbf{supLP120 is the best prior at k=0/k=1}: supLP120 $-$ scratch = \textbf{+5.03 / +2.12 pp} (paired-t p=.0014 / .0153); at k=3 the margin (+2.02 pp) is a trend (p=.067). supLP120 vs supMAE is a wash (36/61 folds, sign p=.20) --- the two priors are indistinguishable on UTD. supMAE is significant-positive at k=0/k=1 (+2.63 / +2.22 pp; p=.0098 / .0255), n.s. at k=3. As on CZU skeleton, \textbf{mae trends positive, not negative} (+2.23/+0.72/+0.59 pp, all n.s.): the reconstruction-hurts signature does not appear on a same-modality small-gap target. Two independent public datasets (CZU skeleton R6, UTD skeleton R6d) agree the pure-supervised prior is at least competitive, and never negative, when the gap is small --- the mirror image of its worst-in-class behavior across the cross-modal gap (R6c).

\textbf{supcon (3 seeds, pooled n=24) is the standout result on this dataset.} supcon $-$ scratch = \textbf{+7.97 / +3.98 / +2.46 pp @ k=0/1/3, all significant (p<.0001 / p=.0002 / p=.028)} --- a clearly larger margin than supLP120's, and supcon beats supLP120 too (sign test 42/60, p=.003, mean$\Delta$ +1.75 pp). Combined with the CZU result (where supcon also edges out supLP120), the two-dataset picture is: \textbf{supcon is the most reliable small-gap performer of the five objectives, ahead of supLP120 on both external datasets}, while supLP120 itself is only conditionally best. The deeper regularity --- "label-supervised, no reconstruction component wins at small gap" --- holds as a class-level claim; which specific label-aware objective wins by how much varies per dataset.

\textbf{Published-baseline comparison.} Same recognizer family as the CZU R6 anchor (statistical moments + CRC-RLS, $\lambda$=1e-4) on the byte-identical seed-42 LOSO splits:

\begin{table*}[!t]
\centering
\caption{R6d --- Second independent same-modality dataset (UTD-MHAD skeleton$\rightarrow$skeleton)}
\small
\begin{tabular}{llllll}
\toprule
k & CRC baseline & scratch & mae & supMAE & supLP120 \\
\midrule
0 & 69.48 & 75.98 & 78.28 & 80.50 & \textbf{83.05} \\
1 & 92.25 & 90.25 & 93.18 & 92.88 & \textbf{93.51} \\
3 & 95.30 & 93.52 & 94.85 & 95.35 & \textbf{96.28} \\
\bottomrule
\end{tabular}
\end{table*}

supLP120 $-$ CRC = \textbf{+13.57 / +1.26 / +0.98 pp} at k=0/1/3, echoing R6's consistent margin over the hand-crafted baseline. One difference from CZU worth noting honestly: on CZU, the \textbf{scratch} learned encoder was statistically level with CRC (R6); on UTD, \textbf{scratch already beats CRC by +6.5 pp at k=0} before any prior is added --- the learned architecture has some edge here even without pretraining, so the prior's marginal contribution on top of scratch is smaller than its margin over CRC. (UTD-MHAD's own published cross-subject accuracy figure is not yet cited here --- TODO, needs a literature figure, not compute.)

\subsubsection{R6e --- Synthesis: the five-setting map}

Five source$\rightarrow$target settings, two public external datasets plus the two internal ones (Xsens position-derived, CZU orientation-only and dual-raw), give a controlled sweep over gap width and target-representation strength:

\begin{table*}[!t]
\centering
\caption{R6e --- Synthesis: the five-setting map}
\small
\begin{tabular}{llllll}
\toprule
\# & Setting (source $\rightarrow$ target) & Gap & Target representation & Best prior ($\Delta$ vs scratch) & Negative-transfer case \\
\midrule
1 & NTU $\rightarrow$ CZU skeleton (R6) & small (same modality) & native skeleton & supLP120 +5.0 pp @k0 (p=.001) & none (mae $\approx$ flat, n.s.) \\
2 & NTU $\rightarrow$ UTD skeleton (R6d) & small (same modality) & native skeleton & supLP120 +5.0 pp @k0 (p=.0014); supcon strongest overall (+7.97pp, p<.0001) & none (mae trends positive, n.s.) \\
3 & NTU $\rightarrow$ Xsens, position-derived quats (R2) & middle (cross-device, skeleton-like target) & strong (mocap-grade quats) & mixed --- supMAE/supLP120/supcon all significant only at k=3 (p$\le$.003), n.s. at k=0/1 & \textbf{mae} (McNemar $-$266/$-$242/$-$99, p<.001 every k) \\
4 & NTU $\rightarrow$ CZU IMU, orientation-only (R6b) & large (cross-modal) & weak (Madgwick quats, yaw drift) & none clearly (supMAE trends > supLP120, 29/43 p=.032) & supLP120 trends worst, n.s. (p=.41--.59) \\
5 & NTU $\rightarrow$ CZU IMU, dual raw (R6c) & large (cross-modal) & strong (raw $\approx$ CRC) & none (\texttt{dual/scratch} itself is best at every k) & \textbf{supLP120} (10/41, p=.0015); \textbf{supcon worse still} (3/42, p<.0001) \\
\bottomrule
\end{tabular}
\end{table*}

No pretraining objective is unconditionally safe across the skeleton$\rightarrow$wearable gap. The clearest pattern: \textbf{mae is the one objective with a real, significant, reproducible negative-transfer signature --- at the middle gap specifically} --- while at small gap no objective is negative and at large gap the pure-label-aware objectives (supLP120, supcon) are the ones that trend or test negative instead. The prior's gap-invariant value is calibration and cold-start data-efficiency (R4).

\textbf{A fifth objective (supcon) sharpens the small/large split into an axis, not two isolated findings.} supcon --- Khosla 2020 SupCon, label-supervised but contrastive rather than softmax --- tracks supLP120's direction at every external setting: it \textbf{wins big at both small-gap datasets} (CZU +5.24/+4.67/+3.41 pp k=0/1/3, all p<.01; UTD +7.97/+3.98/+2.46 pp, all p<.03 --- the single strongest small-gap result in the paper on UTD, clearly ahead of supLP120, p=.003), and it \textbf{trends negative-to-flat at both large-gap settings} (CZU-IMU quat: $-$0.44/$-$1.00/$-$1.04 pp, n.s.; CZU-IMU dual: significantly worse than scratch, p<.0001, the single worst pooled sign-test result in R6c). Because supcon and supLP120 arrive at "label supervision" through different mechanisms (cross-entropy vs contrastive) yet land in the same direction at every gap setting, the operative variable is not "is the objective supervised classification" but \textbf{"does the objective have a reconstruction/regularization component"} --- present in supMAE (the objective that comes closest to never being negative anywhere) and absent from supLP120/supcon alike. This reframes rows 4--5's negative-transfer entries as an instance of a \textit{class} of objectives, not a property specific to softmax classification.

\textit{Gap ordering, measured:} the "middle" placement of NTU$\rightarrow$Xsens is not merely narrative. The encoder-space CKA-per-target check (R3 extension) does \textbf{not} confirm the small/middle/large ordering --- on the supLP120 encoder, mean L0--L2 CKA ranks utd\_skeleton > czu\_skeleton > czu\_imu\_quat > xsens\_v2, putting the claimed-middle target \textit{lowest}, an instance of CKA's own necessary-not-sufficient limitation (R3) rather than a fit we forced. But a raw, encoder-free measurement --- squared-MMD and Frechet distance on a hand-crafted, model-free feature (R3) --- \textbf{does confirm the ordering exactly and monotonically}: czu\_skeleton (MMD\textsuperscript{2}=0.056) < utd\_skeleton (0.096) < xsens\_v2 (0.122) < czu\_imu\_quat (0.205), same ranking on Frechet distance independently. The gap ordering in this table therefore rests on the downstream sign-test/accuracy pattern above, and a measured raw-distributional gap that tracks it directly --- with CKA's alignment-based null read as a separate, orthogonal finding about what encoder alignment does and doesn't predict (R3), not as evidence against the ordering itself.

\subsection{Numbers ledger (traceability)}
\begin{itemize}
  \item R1 within-domain: \texttt{trained\_models/NTU-to-NTU-objective-sanity/}; cross-domain: \texttt{LOSO-fullTrainCalibrate-v2\{,-seed43,-seed44\}/summary.csv} (\texttt{scripts/main\_experiment/loso\_fulltrain\_calibration.py}).
  \item R2 McNemar / R3 CKA / R4a ECE: \texttt{trained\_models/Phase1-analysis/\{mcnemar,cka,cka\_by\_target,ece\}\_results.csv} (\texttt{scripts/main\_experiment/\{dump\_posteriors,cka\_analysis\}.py}).
  \item R3 raw MMD/Frechet: \texttt{trained\_models/RawDomainGap/raw\_domain\_gap.csv} (\texttt{scripts/main\_experiment/raw\_domain\_gap.py}); encoder-space MMD: \texttt{trained\_models/MMD\_DomainGap/mmd\_table.csv} (\texttt{scripts/main\_experiment/mmd\_domain\_gap.py}).
  \item R4b AUC-30/convergence: \texttt{wiki/results/multiseed-loso-v2.md} (\texttt{RESEARCH\_LOG.md} \S{}B).
  \item R4c A2: pooled 3 seeds --- \texttt{trained\_models/A2-subjectScaling\{,-seed43,-seed44\}/N*/summary.csv}, pooled table \texttt{trained\_models/A2-subjectScaling-pooled/\{a2\_pooled\_results,a2\_pooled\_stats\}.csv} (\texttt{scripts/main\_experiment/analyze\_a2\_multiseed.py}).
  \item R5 controller: \texttt{trained\_models/Phase3-controller/robust/\{vocab\_sweep,vocab\_ordering,costmodel\_sweep,costmodel\_summary,frontier,iso\_safety,iso\_safety\_summary\}.csv} + PNGs (\texttt{scripts/controller/controller\_robust.py --vocabs 120 --missions 1000}; all three locks share the same 120 randomly-drawn System Input assignments, no recall-ranked or otherwise non-random vocab used anywhere in the reported numbers).
  \item R6 CZU: \texttt{trained\_models/CZU-skeleton-LOSO\{,-seed43,-seed44\}/summary.csv}; CRC baseline \texttt{.../crc\_baseline/crc\_summary.csv} (\texttt{scripts/external/czu/crc\_baseline.py}).
  \item R6b CZU inertial (cross-modal): \texttt{trained\_models/CZU-IMU-LOSO\{,-seed43,-seed44\}/summary.csv}; CRC \texttt{.../crc\_baseline/crc\_summary.csv} (\texttt{scripts/external/czu/imu\_crc\_baseline.py}); data \texttt{Data\_Processed/czu\_imu\_quats/}.
  \item R6c CZU dual-branch: \texttt{trained\_models/CZU-IMU-DUAL\{,-seed43,-seed44\}/\{raw\_scratch,dual\_<prior>\}/summary.csv} (\texttt{scripts/external/czu/dualbranch.py}); raw data \texttt{Data\_Processed/czu\_imu\_raw/}; reuses CZU-IMU-LOSO splits. Target-richness dial: \texttt{trained\_models/CZU-IMU-DIAL\{,-seed43,-seed44\}/mag20/dual\_<prior>/summary.csv}. Cold-start subject-scaling (T5): \texttt{trained\_models/CZU-DUAL-subjectScaling/N\{0..3\}/dual\_<prior>/summary.csv} (single-seed; 3-seed extension in progress, \texttt{trained\_models/CZU-DUAL-subjectScaling-seed\{43,44\}/}, \texttt{scripts/orchestration/09b\_czu\_cold\_start\_multiseed.sh}).
  \item Pooling for R6/R6b/R6c: \texttt{scripts/external/czu/multiseed\_analyze.py}.
  \item R6d UTD-MHAD skeleton: \texttt{trained\_models/UTD-skeleton-LOSO-seed\{42,43,44\}/summary.csv} (\texttt{scripts/external/utd/crc\_baseline.py} for the CRC row); data \texttt{Data\_Processed/utd\_skeleton\_lrq/}.
  \item OOV / leave-class-out (referenced in \texttt{RESEARCH\_LOG.md}, not a table above): \texttt{trained\_models/LOSO-LeaveClassOutFewShot/summary.csv} (single-seed; 3-seed extension in progress, \texttt{trained\_models/LOSO-LeaveClassOutFewShot-seed\{43,44\}/}, \texttt{scripts/orchestration/05b\_oov\_multiseed.sh}).
\end{itemize}
